\subsection{Non-Injectivity as the Structural Origin of Spinorial Selection}
  \label{subsec:non-injectivity-spinorial}

  The selection of spinorial representations --- characterized here by groups admitting
  a two-dimensional faithful irreducible representation, ultimately converging to $SU(2)$
  --- is not imposed by hand.
  It arises from the interplay between non-injective projection and the bounded-flux
  constraint developed in Chapter~\ref{sec:spectral-admissibility}.

  The key mechanism is the following.
  Non-injective projection introduces a fiber structure over the space of effective
  observables.
  Closed loops in the observable space may fail to lift to closed loops in the relational
  configuration space, generating non-trivial holonomy.
  For holonomy to be non-trivial and physically meaningful, the fiber must support a
  non-abelian structure.

  Within the framework of Chapter~\ref{sec:spectral-admissibility}, the groups capable
  of realising admissible neutral generating sets are classified (Theorem~6.1).
  The classification yields exactly four families:
  $Q_8$, $\mathrm{Dic}_n\ (n\geq 3)$, $2O$, and $2I$.
  All of these are subgroups of $SU(2)$, not of $SO(3)$.

  The exclusion of $SO(3)$ counterparts --- the ordinary dihedral groups, the octahedral
  group $O$, and the icosahedral group $I$ --- follows from the requirement that the
  second-moment tensor $M = \sum_i \vec{u}_i \vec{u}_i^T$ must be compatible with
  the isotropic constraint, and from the binary cover structure that is necessary for
  non-trivial holonomy of the spinorial fiber.
  An $SO(3)$-type generator quotients out the central $\mathbb{Z}_2$ of $SU(2)$, which
  destroys the $4\pi$-periodicity required for spinorial behavior.

  In this sense, the $SU(2)$ thread is not a symmetry postulate.
  It is the unique emergent structure compatible simultaneously with non-injective
  projection, bounded flux, and the non-trivial holonomy required for emergent gauge
  and curvature structure.
